\documentclass[11pt]{article}
\usepackage[margin=1in]{geometry}
\usepackage{amsmath,amssymb,bm,mathtools}
\usepackage{parskip}
\usepackage{titlesec}
\usepackage{hyperref}
\usepackage{booktabs}
\hypersetup{colorlinks=true,linkcolor=blue,urlcolor=blue}

% Headings are set on their own line; body text always starts a new line.
\titleformat{\paragraph}[display]{\normalfont\bfseries}{}{0pt}{}
\titlespacing*{\paragraph}{0pt}{\parskip}{0pt}

\newcommand{\wnom}{\omega_0}
\newcommand{\wsys}{\omega_s}
\newcommand{\wrot}{\omega_r}
\newcommand{\Jmat}{\bm{J}}

\title{SALT Device Models}
\author{}
\date{}

\sloppy
\begin{document}
\maketitle

\section{Scope and conventions}

SALT solves the periodic steady state of a balanced, single-harmonic
transmission network using EMT device models, with the system frequency
$\wsys$ as an unknown. EMT integrates the same device models in the time
domain with the trapezoidal rule.

Every device contributes a residual $F(x)=0$. Newton's method is written in
the stamped form
\begin{equation}
  J(x^{[k]})\,x^{[k+1]} \;=\; J(x^{[k]})\,x^{[k]} \;-\; F(x^{[k]}),
  \label{eq:newton}
\end{equation}
where $x^{[k]}$ is the vector of unknowns at iteration $k$, $F$ is the
stacked vector of device residuals and $J=\partial F/\partial x$ is the
system Jacobian. Written this way, a linear equation contributes its
coefficients to $J$ and its constant to the right-hand side directly.

In the steady-state formulation the network variables are phasors carried
as real and imaginary components $[X_{re},X_{im}]$ at the fundamental, and
machine states in their own rotating frame are DC. Consequently every
$d/dt$ acting on a rotating frame state vanishes, which is noted equation
by equation below.

\subsection{Per-unit bases}
\label{sec:bases}

\begin{tabular}{@{}lll@{}}
\toprule
Quantity & Base & Convention \\
\midrule
Bus voltage        & $V_{\text{base}}^{\text{bus}}$ & line-to-line, RMS, kV \\
Bus current        & $S_{\text{base}}/(V_{\text{base}}^{\text{bus}}\sqrt{3})$ & kA \\
Machine voltage    & $V_{\text{base}}^{\text{gen}}=V_{\text{base}}^{\text{bus}}\!\cdot\!10^3\sqrt{2/3}$ & line-to-neutral, peak, V \\
Machine current    & $I_{\text{base}}^{\text{gen}}=S_{\text{base}}^{\text{gen}}10^6/(V_{\text{base}}^{\text{gen}}\cdot 3/2)$ & peak, A \\
Inverter           & $V_{\text{base}}^{\text{ibr}}$, $I_{\text{base}}^{\text{ibr}}=S^{\text{ibr}}/(V^{\text{ibr}}\sqrt{3})$ & line-to-line, RMS \\
Induction motor    & as machine & line-to-neutral, peak \\
\bottomrule
\end{tabular}

\medskip
Because $V_{\text{base}}^{\text{gen}} = V_{\text{base}}^{\text{bus}}\sqrt{2/3}$
and the conversion from line-to-neutral peak to line-to-line RMS carries
$\sqrt{3/2}$, the machine-to-bus \emph{voltage} scale factor is exactly $1$.
The remaining scale factors
$c_I^{\,\text{gen}\to\text{bus}}$ and $c_V^{\,\text{ibr}\to\text{bus}}$ are:
\begin{equation}
  c_I^{\,\text{gen}\to\text{bus}} = \frac{I_{\text{base}}^{\text{gen}}/\sqrt{2}}{10^3\,I_{\text{base}}^{\text{bus}}},
  \qquad
  c_V^{\,\text{ibr}\to\text{bus}} = \frac{V_{\text{base}}^{\text{ibr}}}{V_{\text{base}}^{\text{bus}}},
\end{equation}
where $c_I^{\,\text{gen}\to\text{bus}}$ converts machine-base current to
bus-base current, $c_V^{\,\text{ibr}\to\text{bus}}$ converts inverter-base
voltage to bus-base voltage, $I_{\text{base}}^{\text{gen}}$ and
$I_{\text{base}}^{\text{bus}}$ are the machine and bus current bases, and
$V_{\text{base}}^{\text{ibr}}$ and $V_{\text{base}}^{\text{bus}}$ are the
inverter and bus voltage bases. Controller parameters (exciter, governor)
are per unit on the machine base and are used without conversion.

\subsection{Sign convention}

Machines and inverters produce power, so their terminal current leaves the
device and enters the network; it is stamped negative into the bus equation.
Loads consume power and are stamped positive.

\subsection{\texorpdfstring{$dq$}{dq}-Transform}

In SALT, the three-phase frequency-domain grid quantities interface with
the $dq$-quantities through a phase shift,\footnote{Because SALT writes
complex quantities as real and imaginary components, the phase shift
$Y = e^{j\theta}X$ is computed as
$\begin{bsmallmatrix}Y_{re}\\ Y_{im}\end{bsmallmatrix}
 = \begin{bsmallmatrix}\cos\theta & -\sin\theta\\
                       \sin\theta & \phantom{-}\cos\theta\end{bsmallmatrix}
   \begin{bsmallmatrix}X_{re}\\ X_{im}\end{bsmallmatrix}$.}
circumventing explicit computation of the $dq$-transform. Different device
models use different conventions of the $dq$-transform, which are
accommodated via linear transformations. The equations for each device are
as follows. Throughout, $\theta$ is the angle of the device's own rotating
frame, $\theta_{a}=\theta$, $\theta_{b}=\theta-\tfrac{2\pi}{3}$ and
$\theta_{c}=\theta+\tfrac{2\pi}{3}$ are the three phase angles, and
$c_V$, $c_I$ are the voltage and current base-conversion factors of
\S\ref{sec:bases}.

\paragraph{Synchronous machine and induction motor.}
The $dq$-transform and its inverse are
\begin{equation}
  \begin{bmatrix}i_d\\ i_q\\ i_0\end{bmatrix}
  = \frac{2}{3}
  \begin{bmatrix}
    \cos\theta_a  & \cos\theta_b  & \cos\theta_c\\
    -\sin\theta_a & -\sin\theta_b & -\sin\theta_c\\
    \tfrac{1}{2}  & \tfrac{1}{2}  & \tfrac{1}{2}
  \end{bmatrix}
  \begin{bmatrix}i_a\\ i_b\\ i_c\end{bmatrix},
  \qquad
  \begin{bmatrix}i_a\\ i_b\\ i_c\end{bmatrix}
  =
  \begin{bmatrix}
    \cos\theta_a & -\sin\theta_a & 1\\
    \cos\theta_b & -\sin\theta_b & 1\\
    \cos\theta_c & -\sin\theta_c & 1
  \end{bmatrix}
  \begin{bmatrix}i_d\\ i_q\\ i_0\end{bmatrix},
\end{equation}
where $i_{a,b,c}$ are the instantaneous phase currents, $i_d,i_q$ are their
direct- and quadrature-axis components in the rotating frame, and $i_0$ is
the zero-sequence component, which vanishes under the balanced assumption.
The SALT equivalent operates on the phase-$a$ phasor:
\begin{equation}
  \begin{bmatrix}I_{\phi a,re}\\ I_{\phi a,im}\end{bmatrix}
  = c_I\,T_{\text{gen}}^{-1}(\theta)\begin{bmatrix}i_d\\ i_q\end{bmatrix},
  \qquad
  T_{\text{gen}}^{-1}(\theta) =
  \begin{bmatrix}\cos\theta & -\sin\theta\\ \sin\theta & \phantom{-}\cos\theta\end{bmatrix},
  \label{eq:Tgen}
\end{equation}
where $I_{\phi a,re}$ and $I_{\phi a,im}$ are the real and imaginary
components of the phase-$a$ current phasor.

\paragraph{Grid-forming inverter.}
The inverter convention negates the quadrature axis,
\begin{equation}
  \begin{bmatrix}i_d\\ i_q\end{bmatrix}
  = \frac{2}{3}
  \begin{bmatrix}
    \cos\theta_a & \cos\theta_b & \cos\theta_c\\
    \sin\theta_a & \sin\theta_b & \sin\theta_c
  \end{bmatrix}
  \begin{bmatrix}i_a\\ i_b\\ i_c\end{bmatrix},
\end{equation}
so the SALT equivalent is the machine phase shift composed with a sign flip
on $q$:
\begin{equation}
  \begin{bmatrix}I_{\phi a,re}\\ I_{\phi a,im}\end{bmatrix}
  = c_I\,T_{\text{ibr}}^{-1}(\theta)\begin{bmatrix}i_{o,d}\\ i_{o,q}\end{bmatrix},
  \qquad
  T_{\text{ibr}}^{-1}(\theta)
  = T_{\text{gen}}^{-1}(\theta)\begin{bmatrix}1&0\\0&-1\end{bmatrix}
  = \begin{bmatrix}\cos\theta & \sin\theta\\ \sin\theta & -\cos\theta\end{bmatrix},
  \label{eq:Tibr}
\end{equation}
where $i_{o,d},i_{o,q}$ are the inverter output-current components.

\paragraph{Composite load.}
The load frame is the system reference, $\theta_{\text{sys}}=0$, so the
phase shift is the identity and the transform disappears entirely:
\begin{equation}
  \begin{bmatrix}v_d^s\\ v_q^s\end{bmatrix}
  = c_V \begin{bmatrix}V_{\phi a,re}\\ V_{\phi a,im}\end{bmatrix},
\end{equation}
where $v_d^s,v_q^s$ are the stator voltage components and
$V_{\phi a,re},V_{\phi a,im}$ the real and imaginary components of the
phase-$a$ bus voltage phasor.

\paragraph{Balanced three-phase output.}
The remaining two phases are enforced algebraically rather than solved: for
a phase-$a$ phasor $I_a$,
\begin{equation}
  \begin{bmatrix}I_b\\ I_c\end{bmatrix}
  = \begin{bmatrix}e^{-j2\pi/3}\\ e^{+j2\pi/3}\end{bmatrix} I_a ,
  \label{eq:balanced}
\end{equation}
which in real and imaginary components reads
\begin{equation}
  \begin{aligned}
    I_{b,re} &= -\tfrac{1}{2}I_{a,re} + \tfrac{\sqrt{3}}{2}I_{a,im}, &\qquad
    I_{b,im} &= -\tfrac{\sqrt{3}}{2}I_{a,re} - \tfrac{1}{2}I_{a,im},\\
    I_{c,re} &= -\tfrac{1}{2}I_{a,re} - \tfrac{\sqrt{3}}{2}I_{a,im}, &\qquad
    I_{c,im} &= \phantom{-}\tfrac{\sqrt{3}}{2}I_{a,re} - \tfrac{1}{2}I_{a,im},
  \end{aligned}
\end{equation}
where $I_{b}$ and $I_{c}$ are the phase-$b$ and phase-$c$ phasors and the
subscripts $re$, $im$ denote their real and imaginary components.

\subsection{Document convention}

Each device model below is stated twice. The first form is the governing
time-domain equation, written with its derivatives intact, as EMT
integrates it. An arrow then points to the steady-state form that SALT
stamps:
\begin{equation*}
  \text{time-domain equation}
  \;\;\xrightarrow{\;\text{SS}\;}\;\;
  \text{steady-state residual stamped by SALT}.
\end{equation*}
Reading left to right therefore shows exactly which terms survive the
steady-state assumption and which vanish. Equations given in one form only
are already identical in both solvers.

\section[Synchronous machine: GENROU]{Synchronous machine:
GENROU\protect\footnote{P.~Kundur, \emph{Power System Stability and
Control}: stator equations $3.120$--$3.134$, electrical torque $3.174$,
machine parameter derivation \S4.1, per-unit bases \S3.4.1 and \S3.4.9,
Park transform $3.60$, initialisation \S3.6.2, current-direction
conventions Fig.~3.9 (generator) and Figs.~7.6--7.7 (load).}}

Round-rotor machine with one field winding ($fd$), one $d$-axis damper
($1d$) and two $q$-axis dampers ($1q$, $2q$).

\subsection{Circuit matrices from nameplate data}

The winding inductances and rotor resistances of the equivalent circuit are
not given directly; they are computed from the nameplate reactances and
open-circuit time constants.
\begin{align}
  L_{ad} &= X_d - X_l, & L_{aq} &= X_q - X_l, & L_\ell &= X_l, & L_0 &= X_l,\\
  L_{ffd} &= \frac{L_{ad}^2}{X_d - X_d'}, & L_{f1d} &= L_{ad}, &
  L_{11d} &= \frac{L_{ad}^2}{X_d - X_d''}, \\
  L_{11q} &= \frac{L_{aq}^2}{X_q - X_q'}, & L_{22q} &= \frac{L_{aq}^2}{X_q - X_d''}.
\end{align}
where $X_d,X_q$ are the synchronous reactances, $X_d',X_q'$ the transient
reactances, $X_d''$ the subtransient reactance, $X_l$ the leakage
reactance, $L_{ad},L_{aq}$ the $d$- and $q$-axis mutual inductances,
$L_\ell$ the leakage inductance, $L_0$ the zero-sequence inductance,
$L_{ffd}$ the field self-inductance, $L_{f1d}$ the field-to-damper mutual
inductance, $L_{11d}$ the $d$-axis damper self-inductance, and
$L_{11q},L_{22q}$ the $q$-axis damper self-inductances.

Rotor resistances follow from the open-circuit time constants. With
\begin{equation}
  L_{fd} = \frac{(X_d'-X_l)L_{ad}}{L_{ad}-(X_d'-X_l)},\quad
  z = X_d''-X_l,\quad
  y = \frac{L_{ad}L_{fd}}{L_{ad}+L_{fd}},\quad
  L_{1d} = \frac{yz}{y-z},
\end{equation}
\begin{equation}
  R_{1d} = \frac{y+L_{1d}}{T_{d0}''},\qquad
  R_{fd} = \frac{L_{ad}+L_{fd}}{T_{d0}'},
\end{equation}
and analogously on the $q$ axis with
$L_{1q}=(X_q'-X_l)L_{aq}/(L_{aq}-(X_q'-X_l))$, $y=L_{aq}L_{1q}/(L_{aq}+L_{1q})$,
$L_{2q}=yz/(y-z)$:
\begin{equation}
  R_{2q} = \frac{y+L_{2q}}{T_{q0}''},\qquad
  R_{1q} = \frac{L_{aq}+L_{1q}}{T_{q0}'} .
\end{equation}
where $T_{d0}',T_{q0}'$ are the open-circuit transient time constants,
$T_{d0}'',T_{q0}''$ the open-circuit subtransient time constants,
$R_{fd}$ the field resistance, $R_{1d}$ the $d$-axis damper resistance,
$R_{1q},R_{2q}$ the $q$-axis damper resistances, $L_{fd},L_{1d},L_{1q},L_{2q}$
the corresponding winding leakage inductances, and $y,z$ intermediate
quantities carrying no physical meaning on their own.

The assembled matrices are
\begin{equation}
  \bm{R} =
  \begin{bmatrix}
    R_a & & & & & & \\
     & R_a & & & & & \\
     & & R_a & & & & \\
     & & & -R_{fd} & & & \\
     & & & & -R_{1d} & & \\
     & & & & & -R_{1q} & \\
     & & & & & & -R_{2q}
  \end{bmatrix},
\end{equation}
where $R_a$ is the armature (stator) resistance, the rotor entries carrying
a negative sign from the equivalent-circuit orientation, and
\begin{equation}
  \bm{L} =
  \begin{bmatrix}
    -(L_{ad}{+}L_\ell) & & & L_{ad} & L_{ad} & & \\
     & -(L_{aq}{+}L_\ell) & & & & L_{aq} & L_{aq}\\
     & & -L_0 & & & &\\
    -L_{ad} & & & L_{ffd} & L_{f1d} & &\\
    -L_{ad} & & & L_{f1d} & L_{11d} & &\\
     & -L_{aq} & & & & L_{11q} & L_{aq}\\
     & -L_{aq} & & & & L_{aq} & L_{22q}
  \end{bmatrix}.
\end{equation}
$\bm{L}_{sp}$ holds the flux terms that multiply rotor speed: it is
$\bm{L}$ with the $d$ and $q$ rows exchanged and signed,
\begin{equation}
  \bm{L}_{sp} =
  \begin{bmatrix}
     & (L_{aq}{+}L_\ell) & & & & -L_{aq} & -L_{aq}\\
    -(L_{ad}{+}L_\ell) & & & L_{ad} & L_{ad} & & \\
     & & & & & & \\
     & & & & & & \\
     & & & & & & \\
     & & & & & & \\
     & & & & & &
  \end{bmatrix}.
\end{equation}

\subsection{Equations}

Let $i=[i_d,i_q,i_{fd}]^\top$, $e=[e_d,e_q,e_{fd}]^\top$, and
$\bm{R},\bm{L}_{sp}$ denote the $(d,q,fd)$ sub-blocks.

\paragraph{Stator.}
\begin{equation}
  e = -\bm{R}\,i + \bm{L}\,\frac{di}{dt} + \bm{L}_{sp}\,i\,\frac{\wrot}{\wnom}
  \;\;\xrightarrow{\;\text{SS}\;}\;\;
  0 = -e + \Big(-\bm{R} + \bm{L}_{sp}\tfrac{\wrot}{\wnom}\Big) i .
\end{equation}
where $i_d,i_q$ are the stator currents in the rotor frame, $i_{fd}$ the
field current, $e_d,e_q$ the corresponding stator voltages, $e_{fd}$ the
field voltage, $\wrot$ the rotor electrical speed and $\wnom$ the nominal
speed.

\paragraph{$dq$-transform, voltage.}
\begin{equation}
  0 = -v_{\phi a} + T_{\text{gen}}^{-1}(\theta)\,c_V\, e_{dq}.
\end{equation}
where $v_{\phi a}$ is the phase-$a$ terminal voltage phasor and
$T_{\text{gen}}^{-1}(\theta)$ is the phase shift of
Equation~\eqref{eq:Tgen}.

\paragraph{Field voltage.}
Supplied by the exciter.

\paragraph{Rotor speed.}
\begin{equation}
  \wrot = \wsys .
\end{equation}
where $\wsys$ is the system frequency. In steady state every machine turns
at the system frequency, so the rotor speed is not an independent unknown.

\paragraph{Swing.}
Electrical torque is derived from the speed terms of the stator equation,
so when $\wrot\neq\wnom$ the mechanical torque is normalised:
\begin{equation}
  T_m\,\frac{\wrot}{\wnom} = e_d i_d + e_q i_q + R_a\big(i_d^2 + i_q^2\big).
  \label{eq:swing}
\end{equation}
where $T_m$ is the mechanical torque supplied by the governor.

In steady state the rotor acceleration is zero, so the swing equation
$\tfrac{2H}{\wnom}\tfrac{d\wrot}{dt} = T_m - T_e$ collapses to
$T_m = T_e$ and is replaced by the conservation of power expressed in
Equation~\eqref{eq:swing}: the mechanical power entering the shaft,
$T_m\wrot/\wnom$, equals the electrical power leaving the machine. The
first two terms on the right-hand side are the electrical power delivered
at the terminals, and the armature term $R_a(i_d^2+i_q^2)$ is the copper
loss dissipated in transferring it.

\paragraph{$dq$-transform, current.}
\begin{equation}
  0 = -I_{\phi a} + T_{\text{gen}}^{-1}(\theta)\,c_I\, i_{dq}.
\end{equation}
where $I_{\phi a}$ is the phase-$a$ terminal current phasor.

\paragraph{Balanced output current.}
Equation~\eqref{eq:balanced}.

\paragraph{Mechanical torque.}
Supplied by the governor.

\section{Exciter: EXDC1}

Five controller states $x_1\ldots x_5$ plus an algebraic terminal-voltage
measurement $v_t$.

\begin{align}
  0 &= -K_e x_1 - T_e\,\dot{x}_1 + x_4
     &&\xrightarrow{\;\text{SS}\;}& 0 &= -K_e x_1 + x_4, \\
  0 &= x_2 + T_r\,\dot{x}_2 - v_t
     &&\xrightarrow{\;\text{SS}\;}& 0 &= x_2 - v_t, \\
  0 &= x_2 + T_c\dot{x}_2 + x_3 + T_b\dot{x}_3 + x_5 + T_c\dot{x}_5 - v_{\text{ref}}
     &&\xrightarrow{\;\text{SS}\;}& 0 &= x_2 + x_3 + x_5 - v_{\text{ref}}, \\
  0 &= -K_a x_3 + x_4 + T_a\,\dot{x}_4
     &&\xrightarrow{\;\text{SS}\;}& 0 &= -K_a x_3 + x_4, \\
  0 &= -K_{f1}\dot{x}_1 + x_5 + T_{f1}\,\dot{x}_5
     &&\xrightarrow{\;\text{SS}\;}& 0 &= x_5 .
\end{align}
where $x_1$ is the exciter output voltage, $x_2$ the measured terminal
voltage, $x_3$ the lead--lag compensator output, $x_4$ the regulator
output, $x_5$ the rate-feedback state, $v_{\text{ref}}$ the voltage
setpoint, $K_e$ and $T_e$ the exciter gain and time constant, $T_r$ the
measurement time constant, $T_b$ and $T_c$ the lead--lag time constants,
$K_a$ and $T_a$ the regulator gain and time constant, and $K_{f1}$ and
$T_{f1}$ the rate-feedback gain and time constant.

Terminal voltage, in machine per unit:
\begin{equation}
  0 = -v_t + \sqrt{\tfrac{2}{3}}\,\Big(v_a^2+v_b^2+v_c^2\Big)^{1/2}.
\end{equation}
where $v_t$ is the terminal-voltage magnitude and $v_a,v_b,v_c$ are the
instantaneous phase voltages.

The exciter drives the machine field-voltage equation through
\begin{equation}
  0 = c_{\text{g}\to\text{e}}\,e_{fd} - x_1\frac{\wrot}{\wnom},
  \qquad c_{\text{g}\to\text{e}} = \frac{L_{ad}}{R_{fd}} .
\end{equation}
where $c_{\text{g}\to\text{e}}$ converts the machine field voltage to the
exciter per-unit base.

\section[Governor: TGOV1D]{Governor:
TGOV1D\protect\footnote{Governor per-unit conventions follow the IEEE PES
report, \url{https://site.ieee.org/fw-pes/files/2013/01/PES_TR1.pdf}.}}

Two states plus an algebraic mechanical power.

\begin{align}
  0 &= x_1 + T_3\dot{x}_1 - x_2 - T_2\dot{x}_2
     &&\xrightarrow{\;\text{SS}\;}& 0 &= x_1 - x_2, \\
  0 &= p_{\text{ref}} - \Big(\tfrac{\wrot}{\wnom}-1\Big) - R x_2 - R T_1 \dot{x}_2
     &&\xrightarrow{\;\text{SS}\;}& 0 &= p_{\text{ref}} - \Big(\tfrac{\wrot}{\wnom}-1\Big) - R x_2, \\
  0 &= x_1 - \Big(\tfrac{\wrot}{\wnom}-1\Big) D_t - p_{\text{mech}} .
\end{align}
where $x_1$ is the turbine power state, $x_2$ the valve position,
$p_{\text{ref}}$ the load reference setpoint, $p_{\text{mech}}$ the
mechanical power delivered to the shaft, $R$ the droop, $T_1$ the governor
time constant, $T_2$ and $T_3$ the turbine lead and lag time constants, and
$D_t$ the turbine damping coefficient.

Electrical and mechanical angles satisfy $\theta = (p/2)\theta_M$, hence
$\wrot = (p/2)\omega_M$, where $p$ is the number of poles, $\theta_M$ the
mechanical rotor angle and $\omega_M$ the mechanical speed. With $p=2$ this
gives $\wrot=\omega_M$ and the mechanical torque coupling into the machine
is
\begin{equation}
  0 = p_{\text{mech}} - \frac{\wrot}{\wnom}T_m .
\end{equation}
where $T_m$ is the mechanical torque appearing in the machine swing
equation.

\section[Grid-forming inverter]{Grid-forming
inverter\protect\footnote{B.~Kenyon et al., ``Open-Source PSCAD
Grid-Following and Grid-Forming Inverters and a Benchmark for Zero-Inertia
Power System Simulations'' (NREL), which is also the source of the default
parameter set; and ``An Accurate Small-Signal Model of Inverter-Dominated
Islanded Microgrids using dq Reference Frame''.}}

Droop-controlled inverter behind an LCL filter, with an inner current loop,
an outer voltage loop and a virtual resistor. All $dq$ quantities are
two-vectors; $\Jmat=\begin{bsmallmatrix}0&-1\\1&0\end{bsmallmatrix}$
represents multiplication by $j$.

\paragraph{Network interface.}
\begin{equation}
  0 = -I_{\phi a} + T_{\text{ibr}}^{-1}(\theta)\,c_I\, i_{o,dq},
\end{equation}
plus Equation~\eqref{eq:balanced}, and
\begin{equation}
  0 = -v_{\phi a} + T_{\text{ibr}}^{-1}(\theta)\,c_V\, v_{g,dq}.
\end{equation}
where $i_{o,dq}$ is the output current at the grid side of the filter,
$v_{g,dq}$ the grid-side voltage, and $T_{\text{ibr}}^{-1}(\theta)$ the
phase shift of Equation~\eqref{eq:Tibr}.

\paragraph{Passive network.}
\begin{align}
  \dot{i}_o &= \tfrac{1}{L_c}\big(-r_c i_o + v_o - v_g\big) + j\omega i_o
    &&\xrightarrow{\;\text{SS}\;}& 0 &= \tfrac{1}{L_c}\big(-r_c i_o + v_o - v_g\big) + \tfrac{\wrot}{\wnom}\Jmat i_o, \\
  \dot{v}_o &= \tfrac{1}{C_f}\big(i_f - i_o\big) + j\omega v_o
    &&\xrightarrow{\;\text{SS}\;}& 0 &= \tfrac{1}{C_f}\big(i_f - i_o\big) + \tfrac{\wrot}{\wnom}\Jmat v_o, \\
  \dot{i}_f &= \tfrac{1}{L_f}\big(-r_f i_f + v_i - v_o\big) + j\omega i_f
    &&\xrightarrow{\;\text{SS}\;}& 0 &= \tfrac{1}{L_f}\big(-r_f i_f + v_i - v_o\big) + \tfrac{\wrot}{\wnom}\Jmat i_f .
\end{align}
where $L_c$ and $r_c$ are the grid-side inductance and resistance, $L_f$
and $r_f$ the converter-side inductance and resistance, $C_f$ the filter
capacitance, $i_f$ the converter-side current, $v_o$ the filter-capacitor
voltage and $v_i$ the converter terminal voltage.

\paragraph{Current controller.}
\begin{align}
  v_i &= k_{C,i}\,\gamma + k_{C,p}\,\dot{\gamma} - j\wrot L_f i_f + G_C v_o
    &&\xrightarrow{\;\text{SS}\;}& v_i &= k_{C,i}\gamma - \tfrac{\wrot}{\wnom}L_f\Jmat i_f + G_C v_o,\\
  \dot{\gamma} &= i_{f}^{\text{ref}} - i_f
    &&\xrightarrow{\;\text{SS}\;}& 0 &= i_{f}^{\text{ref}} - i_f .
\end{align}
where $\gamma$ is the current-loop integrator state, $i_f^{\text{ref}}$ the
current reference from the voltage loop, $k_{C,p}$ and $k_{C,i}$ the
proportional and integral gains, and $G_C$ the voltage feed-forward gain.

\paragraph{Voltage controller.}
\begin{align}
  i_f^{\text{ref}} &= k_{V,i}\,x + k_{V,p}\,\dot{x} - j\wrot C_f v_o + G_V i_o
    &&\xrightarrow{\;\text{SS}\;}& i_f^{\text{ref}} &= k_{V,i}x - \tfrac{\wrot}{\wnom}C_f\Jmat v_o + G_V i_o,\\
  \dot{x} &= v_o^{\text{ref}} - v_o
    &&\xrightarrow{\;\text{SS}\;}& 0 &= v_o^{\text{ref}} - v_o .
\end{align}
where $x$ is the voltage-loop integrator state, $v_o^{\text{ref}}$ the
capacitor-voltage reference, $k_{V,p}$ and $k_{V,i}$ the proportional and
integral gains, and $G_V$ the current feed-forward gain.

\paragraph{Droop.}
With $M_p = \text{droop}_P\,\wnom$ and $M_q = \text{droop}_Q$,
\begin{align}
  \wrot &= \wnom + M_p\big(p_{\text{set}} - p_{\text{avg}}\big), \\
  v_{o,q}^{\text{ref,droop}} &= v_{o,q}^{\text{set}} + M_q\big(q_{\text{set}} - q_{\text{avg}}\big), \\
  \dot{p}_{\text{avg}} &= \omega_{\text{meas}}\big(p - p_{\text{avg}}\big),
    & p &= v_{g,dq}^\top i_{o,dq}
    &&\xrightarrow{\;\text{SS}\;}& 0 &= p - p_{\text{avg}}, \\
  \dot{q}_{\text{avg}} &= \omega_{\text{meas}}\big(q - q_{\text{avg}}\big),
    & q &= v_{g,dq}^\top \Jmat\, i_{o,dq}
    &&\xrightarrow{\;\text{SS}\;}& 0 &= q - q_{\text{avg}}, \\
  \dot{\theta} &= \wrot &&&&\xrightarrow{\;\text{SS}\;}& \wsys &= \wrot .
\end{align}
where $M_p$ and $M_q$ are the active- and reactive-power droop gains,
$p_{\text{set}}$ and $q_{\text{set}}$ the power setpoints, $p$ and $q$ the
instantaneous active and reactive power at the grid side,
$p_{\text{avg}}$ and $q_{\text{avg}}$ their low-pass filtered values,
$\omega_{\text{meas}}$ the measurement filter cutoff,
$v_{o,q}^{\text{set}}$ the nominal quadrature voltage setpoint and
$v_{o,q}^{\text{ref,droop}}$ the droop-adjusted reference.

No factor of $3/2$ appears in $p$ and $q$. It would arise from the
$dq$-transform for peak, line-to-neutral, non-per-unit quantities; with
RMS, line-to-line, per-unit quantities it cancels.

\paragraph{Virtual resistor.}
With $\omega_{VR} = \omega_{VR}^{\text{rad}}/(2\pi)$,
\begin{align}
  \frac{d}{dt}V_o^{VR,\text{avg}} &= \omega_{VR}\big(R_v i_o - V_o^{VR,\text{avg}}\big)
    &&\xrightarrow{\;\text{SS}\;}& 0 &= \omega_{VR}\big(R_v i_o - V_o^{VR,\text{avg}}\big), \\
  V_o^{VR} &= V_o^{VR,\text{avg}} - R_v i_o, \\
  v_o^{\text{ref}} &= \big[0,\;v_{o,q}^{\text{ref,droop}}\big]^\top + V_o^{VR}.
\end{align}
where $R_v$ is the virtual resistance, $\omega_{VR}$ the virtual-resistor
filter cutoff, $V_o^{VR,\text{avg}}$ the filtered virtual-resistor voltage
and $V_o^{VR}$ the voltage correction it applies to the reference.

\section[Composite load: induction motor + ZIP]{Composite load: induction
motor + ZIP\protect\footnote{EPRI, ``Technical Reference on the Composite
Load Model'' (2020), shows that the motor's reactive power cannot be
specified by a power factor and must come from the slip. As such,
initialisation requires its internal Newton--Raphson solve. The default
parameter set comes from the following. The $Z/I/P$ split and the motor
share are from ``Implementation of the WECC Composite Load Model for
Utilities using the Component-Based Modeling Approach'',
\url{https://ieeexplore.ieee.org/stamp/stamp.jsp?tp=&arnumber=7520081}.
Reactive power is taken as pure impedance for industrial loads, following
NERC, ``Reliability Guideline: Developing Load Model Composition Data''
(March 2017),
\url{https://www.nerc.com/globalassets/who-we-are/standing-committees/rstc/lmwg/reliability_guideline_-_load_model_composition_-_2017-02-28.pdf}.
The ZIP low-pass filter time constant is the Simscape default,
\url{https://www.mathworks.com/help/sps/ref/constantpowerloadthreephase.html}.
Motor parameters are Kundur's small industrial motor, \S7.4.3.c; the motor
transform is Kundur $7.23$--$7.27$ and the slip/efficiency relation \S7.7.2.
Supporting literature for the simplified IM${+}$ZIP composition comes from
``Development of Composite Load Models of Power Systems using On-line
Measurement Data'' (2006), ``Standard Load Models for Power Flow and
Dynamic Performance Simulation'' (1995), and ``Load Modeling --- A
Review'' (2018).}}

One induction motor in parallel with one static ZIP load at each load bus.
The motor share of the bus MW is \texttt{powerPercent}.

\subsection{Induction motor}

Let $i_{sr} = [i_{d}^{s}, i_{q}^{s}, i_{d}^{r}, i_{q}^{r}]^\top$ and
$v_{sr} = [v_d^s, v_q^s, 0, 0]^\top$.

\paragraph{Terminal interface.}
\begin{equation}
  0 = -v_{dq}^s + c_V\,v_{\phi a},\qquad
  0 = -\big(i_{dq}^s + i_{dq}^{\text{ZIP}}\big) + c_I\,I_{\phi a},
\end{equation}
plus Equation~\eqref{eq:balanced}, where $v_{dq}^s$ and $i_{dq}^s$ are the
motor stator voltage and current and $i_{dq}^{\text{ZIP}}$ is the current
drawn by the static ZIP load sharing the bus.

\paragraph{Slip and reference angle.}
\begin{equation}
  \omega_{\text{slip}} = \wsys - \wrot, \qquad \theta_{\text{sys}} = 0 .
\end{equation}
where $\omega_{\text{slip}}$ is the slip frequency, $\wrot$ the rotor
speed and $\theta_{\text{sys}}$ the system reference angle, fixed at zero.

\paragraph{Stator and rotor.}
With
\begin{equation}
  \bm{R} = \operatorname{diag}(R_s,R_s,R_r,R_r),\qquad
  \bm{L} = \begin{bmatrix} L_{ss}\bm{I}_2 & L_m\bm{I}_2 \\ L_m\bm{I}_2 & L_{rr}\bm{I}_2\end{bmatrix},
\end{equation}
\begin{equation}
  \bm{\Omega} = \begin{bmatrix}
    0 & -\wsys & 0 & 0\\ \wsys & 0 & 0 & 0\\
    0 & 0 & 0 & -\omega_{\text{slip}}\\ 0 & 0 & \omega_{\text{slip}} & 0
  \end{bmatrix},
\end{equation}
\begin{equation}
  v_{sr} = \bm{R}\,i_{sr} + \frac{d}{dt}\frac{\bm{L}}{\wnom}i_{sr} + \frac{\bm{\Omega}}{\wnom}\bm{L}\,i_{sr}
  \;\;\xrightarrow{\;\text{SS}\;}\;\;
  0 = -v_{sr} + \bm{R}i_{sr} + \bm{\Omega}\,\frac{\bm{L}}{\wnom}i_{sr}.
\end{equation}
where $R_s$ and $R_r$ are the stator and rotor resistances, $L_{ss}$,
$L_{rr}$ and $L_m$ the stator, rotor and magnetising inductances,
$\bm{\Omega}$ the speed-coupling matrix and $i_{sr}$, $v_{sr}$ the stacked
stator and rotor currents and voltages.

The inductances follow from the nameplate reactances as
$L_m = X_m$, $L_{ss} = X_s + X_m$, $L_{rr} = X_r + X_m$, where $X_m$ is the
magnetising reactance and $X_s$, $X_r$ the stator and rotor leakage
reactances; no division by $\omega$ is needed because the reactances are
per unit at $\omega_{pu}=1$.

\paragraph{Swing.}
With
$\bm{L}_{Te} = \bm{L}^\top\!\begin{bsmallmatrix}0&0\\0&0\\0&-1\\1&0\end{bsmallmatrix}\!\begin{bsmallmatrix}0&0&1&0\\0&0&0&1\end{bsmallmatrix}$,
\begin{equation}
  \frac{2H}{\wnom}\frac{d\wrot}{dt} = i_{sr}^\top \bm{L}_{Te}\, i_{sr} - T_0\Big(\frac{\wrot}{\wsys}\Big)^{m}
  \;\;\xrightarrow{\;\text{SS}\;}\;\;
  0 = i_{sr}^\top \bm{L}_{Te}\, i_{sr} - T_0\Big(\frac{\wrot}{\wsys}\Big)^{m}.
\end{equation}
where $H$ is the motor inertia constant, $\bm{L}_{Te}$ the matrix that
extracts electrical torque from the currents, $T_0$ the load torque at
synchronous speed and $m$ the torque--speed exponent of the mechanical
load. As for the synchronous machine, the steady-state form states that the
electrical torque produced balances the mechanical torque absorbed.

\paragraph{Base.}
The motor MVA base is its apparent power at full loading,
\begin{equation}
  S_{\text{base}}^{IM} = \frac{P_{\text{bus}}\cdot\text{powerPercent}}{\mathrm{PF}} ,
\end{equation}
where $P_{\text{bus}}$ is the total active power of the bus,
\texttt{powerPercent} the motor's share of it and $\mathrm{PF}$ the assumed
power factor. Voltage and current bases are derived exactly as for the
synchronous machine. PF only sizes the base; the motor's reactive power is
determined by the slip, not by PF.

\subsection{ZIP}

\begin{align}
  \begin{bmatrix}P\\Q\end{bmatrix}
    &= \begin{bmatrix}P_z\\Q_z\end{bmatrix}V_m^2
     + \begin{bmatrix}P_i\\Q_i\end{bmatrix}V_m
     + \begin{bmatrix}P_p\\Q_p\end{bmatrix}, \\
  V_m &= \sqrt{(v_d^s)^2 + (v_q^s)^2}, \\
  i_d^{\text{ZIP}} &= \frac{P v_d + Q v_q}{V_m^2}, \qquad
  i_q^{\text{ZIP}} = \frac{P v_q - Q v_d}{V_m^2}.
\end{align}
where $P$ and $Q$ are the active and reactive power drawn by the static
load, $V_m$ the terminal-voltage magnitude, $P_z,Q_z$ the
constant-impedance shares, $P_i,Q_i$ the constant-current shares,
$P_p,Q_p$ the constant-power shares and $i_d^{\text{ZIP}},i_q^{\text{ZIP}}$
the resulting load currents.

The $Z$, $I$, $P$ shares are normalised within the ZIP category, so they
sum to one for $P$ and independently for $Q$.

\section{Network elements}

\subsection{Transmission line}

Three-phase nominal-$\pi$ with no mutual coupling. From the case data,
$L = X/\wnom$ and $C = B_{\text{chg}}/(2\wnom)$, where $X$ is the branch
reactance and $B_{\text{chg}}$ the total line-charging susceptance.
\begin{equation}
  Y_{br} = \frac{1}{R + j\omega L},\qquad Y_{sh} = j\omega C,
\end{equation}
\begin{equation}
  I_{\text{from}} = Y_{br}(V_f - V_t) + Y_{sh}V_f,\qquad
  I_{\text{to}}   = -Y_{br}(V_f - V_t) + Y_{sh}V_t .
\end{equation}
where $Y_{br}$ is the series branch admittance, $Y_{sh}$ the shunt
admittance at each end, $R$ the branch resistance, $V_f$ and $V_t$ the
from- and to-bus voltage phasors, and $I_{\text{from}}$, $I_{\text{to}}$
the currents injected at each end.

Because $\omega$ is an unknown, each branch also contributes
$\partial I/\partial\omega$ to the Jacobian column of the frequency
variable:
\begin{equation}
  \frac{\partial I_{br}}{\partial \omega}
  = -\left(\frac{jL}{|Z|^2} + \frac{2\omega L^2 (R - j\omega L)}{|Z|^4}\right)(V_f - V_t),
  \qquad
  \frac{\partial I_{sh}}{\partial \omega} = jC\,V .
\end{equation}
where $Z = R + j\omega L$ is the series branch impedance and $V$ is the
terminal voltage at the end where the shunt is stamped.

\subsection{Transformer}

Series $R$--$L$ with $L = X/\wnom$. When $R\neq 0$ it is stamped as an
admittance $Y=1/(R+j\omega L)$. When $R=0$ there is no admittance, so the
branch current becomes an explicit unknown (modified nodal analysis) and
the branch contributes
\begin{equation}
  0 = -(V_f - V_t) + j\omega L\, I, \qquad
  \frac{\partial}{\partial\omega} = jL\,I .
\end{equation}
where $I$ is the branch current carried as an explicit unknown and $L$ the
leakage inductance.

\subsection{Constant-impedance load}

Selected by the sign of the reactance:
\begin{equation}
  Y =
  \begin{cases}
    \dfrac{1}{R + j\omega L}, & X > 0 \quad (L = X/\wnom), \\[2ex]
    \dfrac{1}{R} + j\omega C, & X < 0 \quad (C = \operatorname{Im}(1/Z)/\wnom), \\[2ex]
    \dfrac{1}{R}, & X = 0 .
  \end{cases}
\end{equation}
where $Y$ is the stamped shunt admittance, $R$ and $X$ the load resistance
and reactance from the case data, $Z = R + jX$, and $L$, $C$ the equivalent
inductance or capacitance the reactance corresponds to.

\subsection{Slack-angle equation}

The system frequency adds one unknown, so one equation is added: the slack
bus angle is pinned by fixing the ratio of the imaginary to the real part
of its phase-$a$ voltage, which is invariant to magnitude,
\begin{equation}
  0 = c\,V_{\text{re}}^{\text{slack}} - V_{\text{im}}^{\text{slack}},
  \qquad c = \frac{\operatorname{Im}(V^{\text{slack}}_{\text{pf}})}{\operatorname{Re}(V^{\text{slack}}_{\text{pf}})}.
  \label{eq:slackangle}
\end{equation}
where $V_{\text{re}}^{\text{slack}}$ and $V_{\text{im}}^{\text{slack}}$ are
the real and imaginary components of the slack-bus phase-$a$ voltage,
$V^{\text{slack}}_{\text{pf}}$ is the slack voltage from the power-flow
solution and $c$ is the tangent of the angle it fixes.

\section{System frequency}

The system frequency in SALT is a variable that is used by device models and
network branch elements, as previously described. It is not explicitly
defined by its own equation, instead emerging naturally from power
mismatches, and equation~(\ref{eq:slackangle}) keeps the system full-rank.
However, during transient simulation, power mismatches are not resolved
explicitly by frequency deviation due to the dynamic flow of power,
requiring EMT to explicitly define the system frequency.

EMT carries the system frequency as an MVA-weighted average of every
rotating machine and inverter speed:
\begin{equation}
  f_{s0} = \frac{\sum_i S_{\text{base},i}\,f_i}{\sum_i S_{\text{base},i}} .
\end{equation}
where $f_{s0}$ is the reported system frequency, $S_{\text{base},i}$ the
MVA base of device $i$ and $f_i$ its instantaneous electrical frequency.

\end{document}
